% Pursuing Stacks (cotes 134-2 to 134-8): archive cote 134-7, site page 52 (PDF page 53), typescript page 494.
% typescript. Draft compared: ps6.tex:2559-2630.
% AI-transcribed by Claude Opus 5.5 (claude-opus-5-5, adaptive thinking, effort high), Message Batches API, batch ps4 (msgbatch_012Aoc1NuskATPsyBeKSYK7J), 2026-11-12; not yet checked by eye.
\scanpage{134-7/52}
\mspage{494}

the dual module
\[\Gamma_k^{\wedge}(M)\simeq(\mathrm{Sym}_k^{*}(M))\hand{^{\vee}}\hand{\simeq}\ \hand{\prod}_j\ (\mathrm{Sym}_k^{j}(M)\hand{^{\vee}}\ \hand{\simeq}\ \Gamma_k^{\,j}(M))\]
\unsure{may|mgy} be viewed as the inverse limit of the duals of those submodules
(compare p.484 (11) ), and topologized accordingly; linear maps (3)
may be interpreted in terms of continuous maps between the dual structures
(or, equivalently, between the corresponding proobjects)
\[(3')\qquad \Gamma_k^{\unsure{\wedge}}(M)\longrightarrow\Gamma_k^{\wedge}(N)\quad ,\]
and compatibility of (3) with multiplication and units is expressed by
compatibility of (3') with``comultiplication'' i.e. diagonal maps, and
with augmentations. On the other hand, maps $W(M)\rightarrow W(N)$ respecting the
``pointed structures'' coming from zero sections, correspond to maps (3')
transforming 1 into 1 , besides the other requirements. Such maps,
it turns out, automatically \xout{transform} induce a map between the submodules
\[(4)\qquad \Gamma_k(M)\longrightarrow\Gamma_k(N)\quad ,\]
which is rather evident indeed, if we \xout{\unsure{recall that}} remind ourselves of \xout{th}
the fact that the submodule $\Gamma_k(M)$ may be viewed as the\inserted{(topological)}dual of $\mathrm{Sym}_k^{*}(M)$\hand{,}
topologized by the powers of the augmentation ideal \deleted{\xout{\unsure{Sym}xx}}
\[\mathrm{Sym}_k^{+}(M)\ =\ \hand{\bigoplus}_{i>0}\ \mathrm{Sym}_k^{i}(M)\quad ,\]
or equivalently of the corresponding adic completion
\[(5)\qquad \mathrm{Sym}_k^{*\wedge}(M)\ =\ \varprojlim_{i}\ \mathrm{Sym}_k^{*}(M)(i)\simeq\ \hand{\prod}_{\unsure{j\geq 0}}\mathrm{Sym}_k^{j}(M)\quad ,\]
and correspondingly for N. The``pointed'' assumption on \deleted{(3)} a map $W(M)\hand{\rightarrow}W(N)$\hand{,}
in terms of the corresponding homomorphism of k-algebras (3), just
translates into compatibility with the augmentations or, equivalently,
with the corresponding ideals $\mathrm{Sym}_k^{+}$ , which implies that it induces a
homomorphism on the corresponding adic rings, and hence by duality a
homomorphism (4) on their duals.

These reminders bring near that working with the\inserted[typed]{(discrete)}k-algebras
$\mathrm{Sym}_k^{*}(M)$, or equivalently with the \uline{topological} coalgebras $\Gamma_k^{\wedge}(M)$, amounts
to working with ``unipotent bundles'' (projective and of finite type),
which are just usual schemes (of a rather particular structure of course);
whereas working with the \uline{topological} k-algebras $\mathrm{Sym}_k^{*\wedge}(M)$, or equivalent
\hand{-}ly with the discrete coalgebras $\Gamma_k(M)$ , amounts to working with \uline{formal}
\uline{schemes} , namely essentially, with the formal completions of the former
along the zero sections.( Of course, the topological objects just consi-

\margin{52}
